% Generated draft: requires mathematical and source review.
\documentclass[11pt]{article}
\usepackage[margin=25mm]{geometry}
\usepackage{fontspec}
\setmainfont{Latin Modern Roman}
\usepackage{amsmath,amssymb,mathtools,unicode-math}
\setmathfont{Latin Modern Math}
\usepackage{enumitem,needspace,xurl,hyperref}
\setlength{\parindent}{0pt}
\setlength{\parskip}{5pt}
\setlength{\emergencystretch}{3em}
\newcommand{\N}{\mathbb{N}}
\newcommand{\Z}{\mathbb{Z}}
\newcommand{\Q}{\mathbb{Q}}
\newcommand{\R}{\mathbb{R}}
\newcommand{\C}{\mathbb{C}}
\newcommand{\F}{\mathbb{F}}
\newcommand{\A}{\mathbb{A}}
\newcommand{\PP}{\mathbb{P}}
\newcommand{\E}{\mathbb{E}}
\newcommand{\eps}{\varepsilon}
\newcommand{\dd}{\,\mathrm{d}}
\begin{document}
\section*{UnsolvedMath ID: 20000430}
\textbf{Source problem code:} catalogue entry
\subsection{Definitions and mathematical statement}
Deligne's question, reported by D. Ramakrishnan, asks whether a Tate class on a variety over a finitely generated characteristic-zero field has algebraic reduction at some prime. The original does not specify smoothness, properness, cohomology, cycle coefficients, reduction model or field of definition of the reduced cycle.

The report's setting is a smooth projective $X/K$, a prime $\ell$, and a class
\[
 c\in H^{2r}_{\mathrm{et}}(X_{\overline K},\Q_\ell(r))
\]
fixed by an open subgroup of $G_K$. This is the Tate-class convention allowing finite field extension. Choose a normal finite-type arithmetic model $S$ and a smooth projective spread $\mathcal X\to S$, with $\ell$ invertible. Good closed points $s$ have finite residue fields $\kappa(s)$; over a field of positive transcendence degree this choice is more data than a rational prime alone. Algebraic means in the image of the cycle map with $\Q_\ell$ coefficients.

For any finite-dimensional space $W$ of these Tate classes, the report takes a common open stabilizer and its normal core, producing a finite Galois quotient of exponent $N$. It claims that every good specialization satisfies pointwise
\[
 \operatorname{Frob}_s^N\operatorname{sp}_s(w)=\operatorname{sp}_s(w)
 \quad\text{for every }w\in W.
\]
The same $N$ works for all good closed points and specialization paths. If the Tate cycle map is surjective for that fiber over the degree-$N$ residue extension $\kappa(s)_N$, namely
\[
 \operatorname{CH}^r(\mathcal X_s\otimes\kappa(s)_N)\otimes\Q_\ell
 \longrightarrow
 H^{2r}_{\mathrm{et}}(\mathcal X_{\bar s},\Q_\ell(r))^{\operatorname{Frob}_s^N=1},
\]
then all of $W$ specializes algebraically simultaneously. In particular one such Tate-valid fiber would answer the existence question for this $W$. Equivalently the specialized classes must vanish in the quotient of the displayed invariant space by the cycle-map image, the report's Tate-defect test.

The report consequently gives all-good-fiber cases: abelian varieties of dimension at most three; K3 surfaces at odd-characteristic good fibers; and cubic fourfolds at good fibers of residue characteristic greater than three. The algebraicity asserted here permits the finite residue extension just specified; it is not automatically a cycle defined over $\kappa(s)$ itself. The report mentions a nonsplit-quadric example showing that distinction can matter, and a Picard-jump obstruction, but supplies neither the example nor the exact obstruction statement. Algebraicity after reduction also does not claim a lift of the cycle back to $X$.

The summary says the general problem remains open. Its specific Tate-theorem inputs, model conventions and certificate proofs are not fully supplied and require independent verification.
\subsection{Short English statement}
Ask whether a characteristic-zero Tate class becomes algebraic at a suitable good reduction. The report supplies a uniform Frobenius exponent for finitely many classes and reduces algebraicity to the Tate conjecture on a fiber over a controlled residue extension. It reports several all-good-fiber cases, while distinguishing geometric algebraicity from a cycle defined over the original residue field.
\subsection{Sources}
\begin{itemize}
\item Local registry: \url{data/unsolvedmath/problems.json}, permanent ID $20000430$, source code \texttt{catalogue entry}; original question and full supplied research summary.
\item ``Problems from the AIM Tate Conjecture Workshop'', question due to P. Deligne and reported by D. Ramakrishnan, item 19 according to supplied verification (extracted as 9): \url{https://aimath.org/WWN/tateconjecture/tateconjecture.pdf}. Canonical location \texttt{aim-algebraic-number-theory-notes.json}, \texttt{notes[81]}; original remarks and literature field empty. No PDF was opened here.
\item No full authoritative references are supplied for the finite-field Tate cases, uniform normal-core certificate, Picard-jump obstruction or nonsplit-quadric example.
\item Supplied machine-generated research attempt 1: partial result, problem open at the run and currently partially solved; rigor metadata ``proved'', novelty confidence low. This non-peer-reviewed aid requires independent verification before citation.
\end{itemize}
\end{document}
